% Figure (Section 7, new at merge): how RBI intervention slows a fall in the rupee.
% Demand shifts right by 1 (A ₹85 -> B ₹90); the RBI sells dollars, supply shifts
% right by 0.6 (-> C ₹87). Same geometry as Figures 4 and 10.
\begin{tikzpicture}[ds, x=1.1cm, y=1.1cm]
\draw[ds axis] (0,0) -- (0,3.7) node[above] {₹ per US dollar};
\draw[ds axis] (0,0) -- (5,0);
\node[below] at (2.5,0) {Quantity of US dollars};
\draw[ds demand] (0.4,3.2) -- (3.6,0.6) node[below] {D\textsubscript{0}};
\draw[ds demand, ds new] (1.4,3.2) -- (4.6,0.6) node[below] {D\textsubscript{1}};
\draw[ds supply] (0.4,0.6) -- (3.6,3.2) node[above] {S\textsubscript{0}};
\draw[ds supply, draw=fxteal] (1.0,0.6) -- (4.2,3.2) node[above right=-2pt, text=fxteal, align=left] {S\textsubscript{1}: RBI\\sells dollars};
\draw[ds guide] (0,1.9) -- (2,1.9) -- (2,0);
\draw[ds guide] (0,2.306) -- (2.5,2.306);
\draw[ds guide] (0,2.062) -- (2.8,2.062) -- (2.8,0);
\draw[ds shift] (2.95,1.2) -- (3.75,1.2);
\draw[ds shift, fxteal] (3.12,2.75) -- (3.58,2.75);
\node[ds eq] at (2,1.9) {}; \node[below=3pt] at (2,1.9) {A};
\node[ds eq=fxorange] at (2.5,2.306) {}; \node[above=3pt, text=fxorange] at (2.5,2.306) {\textbf{B}};
\node[ds eq=fxteal] at (2.8,2.062) {}; \node[below right=1pt, text=fxteal] at (2.8,2.062) {\textbf{C}};
\node[left] at (0,1.86) {85};
\node[left, text=fxteal] at (0,2.08) {\textbf{87}};
\node[left, text=fxorange] at (0,2.33) {\textbf{90}};
\end{tikzpicture}
